\documentclass[11pt]{article}
\usepackage[margin=0.9in]{geometry}
\usepackage{amsmath,amssymb,amsthm}
\usepackage[hidelinks]{hyperref}

\newtheorem*{theorem}{Literature-implied theorem}
\newtheorem*{lemma}{Optimal diagonal paving input}

\newcommand{\sa}{_{\mathrm{sa}}}
\newcommand{\cB}{\mathcal B}
\setlength{\emergencystretch}{3em}

\begin{document}

\begin{center}
{\Large Optimal-Order so-Paving for Type-I MASAs}\\[0.45em]
{\normalsize A scoped literature-implied answer to Conjecture 2.8(2)}\\[0.3em]
{\small Source: Popa--Vaes, arXiv:1412.0631v3.\quad
Support: Ravichandran--Srivastava, arXiv:1706.03737v2.}
\end{center}

\paragraph{Status.}
\texttt{literature\_implied\_answer (partial subcase)}.  The implication below
is not advertised as a theorem in the supporting paper.  It results by replacing
the matrix-paving input in Section 3 of the source paper by the later optimal
multi-paving theorem.

\paragraph{Source question and location.}
Conjecture 2.8(2), page 8 of arXiv:1412.0631v3, asks whether there is a universal
constant $C$ such that
\[
  n_s(x,\varepsilon)\le C\varepsilon^{-2}
\]
for every self-adjoint $x$ and every MASA $A\subset M$.  Here $n_s$ is the
strong-operator paving size from Definition 2.2.  Section 3 proves that every
MASA in a type-I von Neumann algebra with separable predual is so-pavable, with
the then-available estimate $O(\varepsilon^{-4})$.

\paragraph{Supporting theorem and location.}
Ravichandran--Srivastava, Theorem 1 on page 1 of arXiv:1706.03737v2, states
that $k$ zero-diagonal Hermitian contractions admit a common one-sided
$\eta$-paving with at most $18k\eta^{-2}$ blocks.  Applied to $T$ and $-T$,
this gives a two-sided paving
\[
  \max_j\|P_jTP_j\|<\eta\|T\|,
  \qquad r\le 36\eta^{-2}.                                      \tag{1}
\]

\begin{theorem}
Let $M$ be a type-I von Neumann algebra with separable predual and let
$A\subset M$ be an arbitrary MASA.  For every $x=x^*\in M$ and
$0<\varepsilon<1$,
\[
  n_s(A\subset M;x,\varepsilon)
  \le \left\lceil 144\varepsilon^{-2}\right\rceil .             \tag{2}
\]
Consequently Conjecture 2.8(2) has the optimal order $\varepsilon^{-2}$ on
the entire type-I class.
\end{theorem}

\paragraph{Proof.}
We follow exactly the type-I reduction and lifting argument of Section 3 in
the source paper, changing only its finite-matrix estimate.

First consider a discrete diagonal field
\[
  \mathcal D\,\overline\otimes L^\infty(X)
  \subset \cB(\ell^2 I)\,\overline\otimes L^\infty(X).
\]
For finite $I$, (1) applies pointwise.  For arbitrary countable $I$, apply the
finite theorem to increasing finite corners and pass to a coloring in the
compact product space $\{1,\ldots,r\}^{I}$.  For a measurable field, the
successful colorings form a Borel subset of that compact space times $X$;
the measurable-selection step in Lemma 3.1 of the source paper therefore
chooses them measurably, with the same $r$.

For arbitrary $x=x^*$ in such a field, put $a=E_A(x)$ and $T=x-a$.  Then
$E_A(T)=0$, $\|a\|\le\|x\|$, and $\|T\|\le2\|x\|$.  Use (1) with
$\eta=\varepsilon/2$.  The resulting number of blocks is at most
$36(\varepsilon/2)^{-2}=144\varepsilon^{-2}$ and
\[
  \left\|\sum_j p_j(x-a)p_j\right\|
  =\max_j\|p_jTp_j\|
  \le \frac{\varepsilon}{2}\|T\|
  \le \varepsilon\|x\|.                                       \tag{3}
\]
Thus (2) even holds as a norm paving on every discrete expected component.

It remains to recall the continuous type-I component used in Proposition 3.2
of the source paper:
\[
 L^\infty([0,1])\,\overline\otimes L^\infty(X)
 \subset
 \cB(L^2[0,1])\,\overline\otimes L^\infty(X).                   \tag{4}
\]
Fix a strong neighborhood $\mathcal V$ of zero.  Let $q_m$ project onto the
step functions that are constant on the $m$ equal subintervals of $[0,1]$;
choose $m$ so that $1-q_m\otimes1\in\mathcal V$.  If
$V_m:\mathbb C^m\to L^2[0,1]$ is the corresponding step-function isometry,
compress $x$ to
\[
 y=(V_m^*\otimes1)x(V_m\otimes1).
\]
As in the source proof, let $b$ be the diagonal of $y$ and lift $b$ to an
element $a\in A$ that is constant on each interval.  Then
$\|a\|\le\|x\|$, the matrix field $T=y-b$ is zero diagonal, and
$\|T\|\le2\|x\|$.

Apply (1), again with $\eta=\varepsilon/2$, and choose the fiberwise
partitions measurably.  Lift their diagonal matrix projections to a partition
$p_1,\ldots,p_r\in A$.  The computation in Proposition 3.2 gives
\[
 \left\|(q_m\otimes1)
   \left(\sum_jp_jxp_j-a\right)(q_m\otimes1)\right\|
 =\left\|\sum_j f_j(y-b)f_j\right\|
 \le\varepsilon\|x\|.                                         \tag{5}
\]
The number $r\le\lceil144\varepsilon^{-2}\rceil$ is independent of
$\mathcal V$, which is precisely the required so-paving estimate.  The
standard direct-sum decomposition of a type-I algebra and padding partitions
with zero projections completes the proof. \qed

\paragraph{Sharpness and scope.}
The order cannot be improved uniformly over type-I MASAs: the diagonal MASA
in $\cB(\ell^2\mathbb N)$ already has the lower bound
$\Omega(\varepsilon^{-2})$, as recalled in both papers.  The theorem does not
settle so-paving for a general non-type-I MASA and does not establish the
universal constant requested by Conjecture 2.8 in that generality.

\paragraph{Search status.}
A bounded search on 2026-08-13 used the exact source title and the phrases
``so-paving MASA'', ``Popa Vaes so-paving conjecture'', and ``strong operator
paving MASA''.  It found no claimed solution of the arbitrary-MASA conjecture.
The supporting authors cite Popa--Vaes while discussing optimal multi-paving,
but do not explicitly state (2); the relation is an agent-identified
implication.

\paragraph{References.}
\begin{itemize}
\item S. Popa and S. Vaes, \emph{Paving over arbitrary MASAs in von Neumann
algebras}, arXiv:1412.0631v3 (2015).
\item M. Ravichandran and N. Srivastava, \emph{Asymptotically Optimal
Multi-Paving}, arXiv:1706.03737v2 (2017).
\item P. Casazza, D. Edidin, D. Kalra, and V. Paulsen,
\emph{Projections and the Kadison--Singer problem}, Oper. Matrices 1 (2007),
391--408.
\end{itemize}

\end{document}
